\section{Weak Coupling and Renormalization}
\label{sec:weak_coupling}

In the limit $\beta \to \infty$ (weak coupling), the lattice theory must define the continuum physics. This regime is governed by Asymptotic Freedom.

\subsection{Perturbative Expansion}
Standard lattice perturbation theory allows for the expansion of expectation values in powers of $g_0^2 = 2N/\beta$.
\begin{equation}
\langle \mathcal{O} \rangle = \sum_{k=0}^\infty c_k (g_0^2)^k
\end{equation}
While this series is known to be asymptotic (divergent), it provides the definition of the running coupling via the Callan-Symanzik equation.

\subsection{Renormalization Group Equation}
The beta function $\beta(g)$ is given at one-loop order by:
\begin{equation}
\beta(g) = -\frac{11N}{48\pi^2} g^3 + O(g^5)
\end{equation}
The negative sign implies asymptotic freedom.

\subsection{Scale Setting}
The physical mass scale $m_{phys}$ is related to the lattice spacing $a$ by:
\begin{equation}
m_{phys} a \sim \exp\left( -\int^g \frac{dg'}{\beta(g')} \right)
\end{equation}
Rigorous control of this limit requires establishing non-perturbative bounds on the correlation length $\xi(g)$.

\subsection{Open Problem: Large Field Fluctuations}
The primary obstacle to a rigorous construction in this regime is the control of large field fluctuations (renormalons and non-perturbative effects) which are not captured by the perturbative expansion. The "Probability of Large Fields" must be bounded uniformly as the lattice spacing goes to zero. This remains the central analytic challenge of the Millennium Problem.
